\scanpage{327}
\mspage{66}
% archive sheet number: 326
% ms. 1–6
sont dans $X$. On notera que pour qu'un chemin \inserted{$c$} et son opposé $c^{\unsure{\circ|o}}$ soient tous les deux des \uline{$\Sigma$-chemins}, il \unsure{f. et s.} que \uline{toutes} les flèches du chemin $c$ soient dans $\Sigma$ (on \unsure{dira} donc que $c$ est un \uline{$\Sigma$-chemin strict}), puisque les flèches rétrogrades de $c^\circ$ sont les flèches \deleted{\unsure{strictes}} directes de $c$.

% ms. 7–8
Pour que les sous-catégories pleines $\underline{\mathrm{Ch}}_{\tau,\Sigma}(X)$ soient stables pour les foncteurs \unsure{images inverses|image inverse}
% ms. 9
\[f^*:\underline{\mathrm{Ch}}_{\tau'\deleted{\unsure{\Sigma}}}(X)\unsure{\mathord{-}|\longrightarrow}\underline{\mathrm{Ch}}_\tau(X)\]
% ms. 10
associés aux flèches
% ms. 11
\[f:\tau\longrightarrow\tau'\]
% ms. 12–15
dans $\underline{\unsure{\mathrm{Ty}|\mathrm{Tz}}}$, ou simplement \deleted{\unsure{pour}} qu' \deleted{\unsure{elles soient st}}\inserted{il y ait} stabilité pour \deleted{les deux} foncteurs $f^*$ \deleted{qui} \unsure{associés|associé}, \unsure{au} $f:\Delta_-\longrightarrow1$ \deleted{\unsure{(avec $\tau\to1$)}}\inserted{et \deleted{ainsi que pour} $\Delta_-\longrightarrow\Delta_-^2$}, il f. et s. que $\Sigma$ \inserted{soit stable par composition et} contienne toutes les flèches identiques \unsure{:|;}
% ms. 16–17
\[\left\{\begin{aligned}x\in\deleted{\unsure{\mathrm{Ob}}}X&\Longrightarrow\mathrm{id}_x\in\Sigma.\\
u,v\in\Sigma,\ u,v\ \text{\unsure{composables}}&\Longrightarrow uv\in\Sigma\end{aligned}\right.\]
% ms. 18–26
En effet, les flèches dans un chemin $f^*(c')\in\mathrm{Ch}_\tau(X)$ sont soit \deleted{elles} des flèches identiques (qui correspondent aux \deleted{\unsure{couples}} flèches de transition de $I_\tau$ dégénérées pour $f$), soit des flèches parmi celles de $\tau$, ou des composées d'une suite de flèches \deleted{\unsure{compos}}\inserted{consécutives} composables de $\deleted{\unsure{\Sigma}}\tau'$ (qui correspondent aux \deleted{\unsure{couples}}\inserted{flèches} de transition dans $\unsure{I_{\tau'}|\bar I_{\tau'}}$ incluses dans la portée d'une flèche de transition de $I_\tau$), avec conservation des \unsure{signatures}.
